%%%%%%%%%%%%%%%%%%%%%%%%%%%%%%%%%%%%%%%%%
% Medium Length Professional CV
% LaTeX Template
% Version 2.0 (8/5/13)
%
% This template has been downloaded from:
% http://www.LaTeXTemplates.com
%
% Original author:
% Rishi Shah 
%
% Important note:
% This template requires the resume.cls file to be in the same directory as the
% .tex file. The resume.cls file provides the resume style used for structuring the
% document.
%
% Modified by Sterling G. Baird
%
%%%%%%%%%%%%%%%%%%%%%%%%%%%%%%%%%%%%%%%%%
%
% Teaching Interests
% - https://mitcommlab.mit.edu/eecs/commkit/faculty-application-teaching-statement/
% - https://careerservices.upenn.edu/application-materials-for-the-faculty-job-search/teaching-philosophies-for-faculty-job-applications/
% - 
% Research Interests
% - https://mitcommlab.mit.edu/eecs/commkit/faculty-application-research-statement/
% - https://mitcommlab.mit.edu/eecs/commkit/faculty-application
% - https://gradschool.cornell.edu/career-and-professional-development/pathways-to-success/prepare-for-your-career/take-action/research-statement/
%
% CV
% - https://mitcommlab.mit.edu/eecs/commkit/faculty-application-cv/


%----------------------------------------------------------------------------------------
%	PACKAGES AND OTHER DOCUMENT CONFIGURATIONS
%----------------------------------------------------------------------------------------

\documentclass{resume} % Use the custom resume.cls style

\usepackage[left=0.75in,top=0.6in,right=0.75in,bottom=0.6in]{geometry} % Document margins
\usepackage{hyperref}
\usepackage{needspace}
\usepackage{fontawesome5}
\usepackage{academicons}

\usepackage{xcolor}
\hypersetup{
    colorlinks,
    linkcolor={red!50!black},
    citecolor={blue!50!black},
    urlcolor={blue!80!black}
}

\usepackage{doi}
\usepackage{array}
\usepackage [autostyle, english = american]{csquotes}
\MakeOuterQuote{"}

\usepackage[backend=biber, natbib=true, style=ieee, sorting=none, doi=true, isbn=false, url=true]{biblatex}

% Redefine the urldate bibmacro to do nothing
\renewbibmacro{urldate}{%
  % Empty definition
}

\makeatletter
% auxiliary bibfile
\def\hlblx@bibfile@name{\jobname -boldnames.bib}
\newwrite\hlblx@bibfile
\immediate\openout\hlblx@bibfile=\hlblx@bibfile@name

\newcounter{hlblx@name}
\setcounter{hlblx@name}{0}

\newcommand*{\hlblx@writenametobib}[1]{%
  \stepcounter{hlblx@name}%
  \edef\hlblx@tmp@nocite{%
    \noexpand\AfterPreamble{%
      \noexpand\setbox0\noexpand\vbox{%
        \noexpand\hlblx@getmethehash{hlblx@name@\the\value{hlblx@name}}}}%
  }%
  \hlblx@tmp@nocite
  \immediate\write\hlblx@bibfile{%
    @misc{hlblx@name@\the\value{hlblx@name}, author = {\unexpanded{#1}}, %
          options = {dataonly=true},}%
  }%
}

\AtEndDocument{%
  \closeout\hlblx@bibfile}

\addbibresource{\hlblx@bibfile@name}

\newcommand*{\hlbxl@boldhashes}{}
\DeclareNameFormat{hlblx@hashextract}{%
  \xifinlist{\thefield{hash}}{\hlbxl@boldhashes}
    {}
    {\listxadd{\hlbxl@boldhashes}{\thefield{fullhash}}}}

\DeclareCiteCommand{\hlblx@getmethehash}
  {}
  {\printnames[hlblx@hashextract][1-999]{author}}
  {}
  {}

% user-level macros
\newcommand*{\addboldnames}{\forcsvlist\hlblx@writenametobib}
\newcommand*{\resetboldnames}{\def\hlbxl@boldhashes{}}


\newcommand*{\mkboldifhashinlist}[1]{%
  \xifinlist{\thefield{hash}}{\hlbxl@boldhashes}
    {\mkbibbold{#1}}
    {#1}}
\makeatother

\DeclareNameWrapperFormat{given-family:hash:bold}{%
  \renewcommand*{\mkbibcompletename}{\mkboldifhashinlist}%
  #1}

\DeclareCiteCommand{\printpublication}
  {\usebibmacro{prenote}}
  {\usedriver
     {\DeclareNameAlias{sortname}{default}%
      \DeclareNameWrapperAlias{default}{given-family:hash:bold}%
      \defcounter{maxnames}{12}%
      \defcounter{minnames}{3}}
     {\thefield{entrytype}}}
  {\multicitedelim}
  {\usebibmacro{postnote}}

\addboldnames{{Baird, Sterling G.}}
\addboldnames{{Baird, Sterling Gregory}}
\addboldnames{{Baird, Sterling}}
\addboldnames{{Baird*, Sterling}}
\addbibresource{mypublications.bib}
\addbibresource{custom-publications.bib}

% \bibliography{mypublications.bib}

\newcommand{\tab}[1]{\hspace{.2667\textwidth}\rlap{#1}}
\newcommand{\itab}[1]{\hspace{0em}\rlap{#1}}

\usepackage{longtable}

\name{Sterling G. Baird} % Your name
\address{4138 Devonshire Dr. Provo UT 84604} % Your address
%\address{123 Pleasant Lane \\ City, State 12345} % Your secondary addess (optional)
\address{+1(720)854-5562 \\ \href{mailto:sterling.baird@byu.edu}{sterling.baird@byu.edu} \\ \href{https://www.linkedin.com/in/sterling-baird/}{\faLinkedin\ } 
    % \href{https://x.com/SterlingBaird1}{\faTwitter\ } 
    \href{https://github.com/vertical-cloud-lab}{{\small\faFlask}\ }
    \href{https://github.com/sgbaird}{\faGithub\ } \href{https://scholar.google.com/citations?hl=en&user=UACmnBgAAAAJ}{\aiGoogleScholar\ }
    % \href{https://github.com/vertical-cloud-lab}{\faUsers\ }
    } % Your phone number and email
    
% \address{\href{https://www.linkedin.com/in/sterling-baird/}{LinkedIn} \\ \href{https://x.com/SterlingBaird1}{X} \\ \href{https://github.com/sgbaird}{GitHub}} % LinkedIn, X, GitHub

\begin{document}
% \bibliographystyle{plainurl}
% \nobibliography{my-publications.bib}
%----------------------------------------------------------------------------------------
%	EDUCATION SECTION
%----------------------------------------------------------------------------------------

\begin{rSection}{Education}

{\bf University of Utah (UoU), Salt Lake City, UT}
\\ Ph.D. Materials Science and Engineering \hfill {\em August 2020 - August 2023}
\\ Department of Materials Science and Engineering %\hfill {\em Cumulative GPA: 3.9}

{\bf Brigham Young University (BYU), Provo, UT} 
\\ M.S. Mechanical Engineering \hfill {\em August 2018 - April 2021}
\\ Department of Mechanical Engineering %\hfill {\em Cumulative GPA: 3.6} \linebreak %3.59
\\ B.S. Applied Physics \hfill {\em June 2011 - August 2018}
\\ Department of Physics and Astronomy %\hfill {\em  Major GPA: 3.2} %3.18

{\bf Colorado School of Mines (CSM), Golden, CO} \hfill {\em August 2015 - December 2015}
%\\ Non-degree seeking \hfill {\em Cumulative GPA: 4.0}

%Minor in Linguistics \smallskip \\
%Member of Eta Kappa Nu \\
%Member of Upsilon Pi Epsilon \\


\end{rSection}

%\begin{rSection}{Primary Objective}
%To better standard of living through materials innovation.
%\end{rSection}
%--------------------------------------------------------------------------------
%    Projects And Seminars
%-----------------------------------------------------------------------------------------------
% emphasize skills related to automation
\Needspace{10\baselineskip} % Adjust the space needed accordingly
\begin{rSection}{Research Experience}
\begin{rSubsection}{Brigham Young University (Mech. Eng.), Assistant Professor}{Dec 2025 - present}{} \\
\item Establishing the \href{https://github.com/vertical-cloud-lab}{Vertical Cloud Lab} for autonomous materials discovery and advanced manufacturing
\item Focusing on autonomous alloy discovery via additive manufacturing for aerospace applications
\item Developing cloud-accessible experimentation platforms with remote programmatic access
\end{rSubsection}

\begin{rSubsection}{Acceleration Consortium, Principal Investigator, Staff Scientist}{Mar 2025 - Dec 2025}{} \\
\item Implemented adaptive experimentation strategies within the AC's core physical science facilities
\item Led research initiatives within the AI \& Automation Lab
\item Advanced autonomous laboratory methodologies for accelerated materials discovery
\end{rSubsection}

\begin{rSubsection}{Meta, Technical Consultant}{Nov 2024 - Jan 2025}{} \\
\item Creating docs and refining code for the v1.0 release of Ax Platform, a Bayesian optimization package
\item Providing insights and perspectives on Bayesian optimization applied to the physical sciences
\end{rSubsection}

\begin{rSubsection}{Acceleration Consortium, Director of Training and Programs}{Jun 2023 - Mar 2025}{} \\
\item Developed \href{https://honegumi.readthedocs.io/en/latest/}{\texttt{Honegumi}}, an interactive script generator for materials-relevant Bayesian optimization
% \item Leveraged AI and automation for accelerated stress testing in an electrolysis cell for fuel cell applications
\item Developed \href{https://huggingface.co/collections/AccelerationConsortium/optimization-benchmarks-66a44daf10de1a0335f28826}{Bayesian optimization benchmarks/web apps} for chem/mat.sci. with Merck KGaA partner
\item Leveraged AI and automation for accelerated stress testing in an electrolysis cell for fuel cell applications
\item Leveraged AI and automation for autonomous battery slurry optimization of viscosity and aggregates
\item Developed dozens of automation hardware APIs and tools (e.g., solid dosing, sample transfer)
\item \href{https://github.com/ACC-HelloWorld/1-running-the-demo/}{Incorporated remote-access equipment} into autograded course workflows and assignments
\item Built the \href{https://ac-training-lab.readthedocs.io/}{AC Dev Lab} from scratch, including \$600k of equipment (500 line items) over two years
\item Conducted a feasibility study for the incorporation of vertical lift modules into wetlab environments % a "vertical lab" using a vertical lift module and a mobile manipulator
\end{rSubsection}

\begin{rSubsection}{UoU Sparks Group PhD Research Assistant}{Aug 2020 - Aug 2023}{} \\
\item Developed alloy-focused microstructure informatics software for Fehrmann Materials GmbH
\item \href{https://doi.org/10.21105/joss.04528}{Applied diffusion generative models} for novel crystal structure generation
\item \href{https://www.nsf.gov/awardsearch/showAward?AWD_ID=2334411&HistoricalAwards=false}{Co-wrote} a successful 200k USD grant for scientific literature data extraction via LLMs
\item Developed modular control, communication, data, and optimization solutions for self-driving labs
\item Designed an off-the-shelf, low-cost "Hello World" for self-driving labs (350 kits in use)
% \item Designed an off-the-shelf, low-cost flow reactor demonstration as a "Hello, World!" extension
\item Co-wrote the structural materials section of a \href{10.1021/acs.chemrev.4c00055}{Chem. Rev. review paper} focused on self-driving labs
% \item Wrote and managed a \href{https://doi.org/10.1039/D3DD00223C}{review paper for low-cost self-driving labs}
% \item Created a tool to convert between crystal structure and PNG images for materials generative models
\item \href{https://doi.org/10.1016/j.commatsci.2023.112134}{Applied constrained, high-dimensional Bayesian optimization} for solid rocket propellant simulations
\item Created an \href{https://doi.org/10.1039/D1DD00028D}{adaptive design framework} for high-performance, chemically unique materials discovery
% \item Created a brute-force computational geometry tool for atomic planar densities in arbitrary materials
\item Collected and analyzed a \href{https://doi.org/10.1016/B978-0-12-823144-9.00079-0}{corpus of 50 validated materials discoveries} (computational and experimental)
\end{rSubsection}

\begin{rSubsection}{BYU Materials Innovation Lab Research Assistant}{Aug 2018 - Aug 2021}{} \\
\item Developed a state-of-the-art package for grain boundary property prediction (50\% improvement)
\item Numerically validated and gained new insights to a 20-year assumption about grain boundaries % \cite{baird_quantitative_2021}
\item Designed a semi-automated, ultra-high temp./ultra-high purity induction furnace (2500 °C)
\item Designed a semi-automated, 3D printed, two-cell electrochemical setup
\item Performed image segmentation and multi-modal stitching on microscope and synchrotron data
\item Used image-based techniques (e.g., mutual information) to infer subsurface grain boundary tilt angles % \cite{fullwood_determining_2021}
\item Used Bayesian inference to infer local material properties from global network/graph measurements % \cite{kurniawan_grain_2019,snow_grain_2021,johnson_inference_2021}
\end{rSubsection}

\begin{rSubsection}{BYU Condensed Matter Research Assistant}{Jun 2015 - Aug 2018}{} \\
% \item Synthesized structurally modified carbon nanotube samples for Si-anode Li-ion batteries
\item Performed 100's of experiments using structurally modified carbon nanotubes in lithium-sulfur batteries
\item Reviewed several hundred next-generation battery articles to create a visual database tool
\end{rSubsection}
\begin{rSubsection}{NASA Langley Research Center Intern}{Jan 2017 - May 2017}{} \\
\item Fabricated/tested dozens of lithium-oxygen batteries using structurally modified graphene electrodes
\item Performed Raman laser and FTIR spectroscopic studies on boron nitride nanotube materials
\end{rSubsection}
\begin{rSubsection}{CSM Research Assistant}{Sep 2015 - Oct 2015}{} \\
\item Characterized thin-film solar cell materials via capacitance-voltage and X-ray diffraction measurements
\end{rSubsection}
\end{rSection}

%-----------------------------------
% LEADERSHIP EXPERIENCE
%-----------------------------------
% \pagebreak
\Needspace{10\baselineskip} % Adjust the space needed accordingly
\begin{rSection}{Leadership Experience}

% \begin{rSubsection}{Brigham Young University (Mech. Eng.), Assistant Professor}{Dec 2025 - present}{} \\
% \item Building a research group \href{https://github.com/vertical-cloud-lab}{Vertical Cloud Lab}
% \end{rSubsection}

\begin{rSubsection}{Brigham Young University (Mech. Eng.), Assistant Professor}{Dec 2025 - present}{} \\
\item Supervising 8 undergraduate research assistants from mechanical engineering and applied math
\end{rSubsection}
\begin{rSubsection}{Acceleration Consortium, Principal Investigator and Staff Scientist}{Mar 2025 - Dec 2025}{} \\
\item Supervised 8 students as lab liaisons for research projects (e.g., liquid-liquid extraction, drug discovery)
\end{rSubsection}
\begin{rSubsection}{Acceleration Consortium, Director of Training and Programs}{Jun 2023 - Mar 2025}{} \\
\item Supervised 12 staff on training \& research projects (e.g., course dev., electrolysis, Bayesian opt.)
\item Mentored two visiting Ph.D. MSE researchers in the \href{https://ac-training-lab.readthedocs.io/}{AC Dev Lab} (cumulative 10 weeks, full-time)
\item Led the development of course content, workshops, and an in-person training facility
\item Solo-organized the \href{https://ac-bo-hackathon.github.io/}{AC BO Hackathon} comprising 120 participants, 62 organizations, and 14 countries
\item Spearheaded the \href{https://accelerated-discovery.org}{Accelerated Discovery forum} (31k page views, 270 users, \& 72 discussions in 8 months)
\item Organized scientific technical content for Accelerate 2023 (60 speakers, 100 posters, 350 attendees)
\item Organized nine \href{https://youtu.be/IVaWl2tL06c?si=DYrhZP6uvK5srQ58}{workshops} for data science, robotics, and software development topics
\item Leading a community-curated \href{https://github.com/AccelerationConsortium/awesome-self-driving-labs}{set of resources} for autonomous experiments (126 GitHub stars)
% \item Delivering industry workshops and technical advising to industry partners
\end{rSubsection}
\begin{rSubsection}{Review of low-cost self-driving laboratories}{July 2022 - Feb 2024}{} \NLS
\item Led a \href{https://doi.org/10.1039/D3DD00223C}{review paper for low-cost self-driving labs} with 12 co-authors
\end{rSubsection}
\begin{rSubsection}{Generative Modeling for Crystal Structures}{Jun 2022 - Jun 2023}{} \\
\item Led a team of four scientists/engineers in developing a tool for generative modeling of crystal structures
% \item Supervision, project administration, conceptualization, and division of labor
\end{rSubsection}
\begin{rSubsection}{Data-driven materials discovery and synthesis}{Sep 2020 - Feb 2021}{} \NLS
\item Led a team of three data scientists in creating a corpus of validated materials discovery articles
\item Resulted in accepted publication in \textit{Comprehensive Inorganic Chemistry III}
\end{rSubsection}
\begin{rSubsection}{Five degree-of-freedom hydrogen permeation measurements}{Aug 2020 - Feb 2021}{} \\
\item Led a team of three eng. students towards creating an experimental grain boundary $\mathrm{H}_2$ diffusion model
\end{rSubsection}
\begin{rSubsection}{Mentorships}{Aug 2020 - present}{} \\
\item \textit{University of Utah}: Mentored grad, undergrad, and secondary students (accepted to Stanford, Yale) \smallskip
\item \textit{Brigham Young University}: Mentored two students (wetlab, software skills, career advisement)
\end{rSubsection}
\end{rSection}

%-----------------------------
% TEACHING EXPERIENCE
%-----------------------------
\Needspace{10\baselineskip} % Adjust the space needed accordingly
\begin{rSection}{Teaching Experience}

\begin{rSubsection}{Brigham Young University (Mech. Eng.), Assistant Professor}{Dec 2025 - present}{} \\
\item Instructor for Introduction to Materials Science
\end{rSubsection}

\begin{rSubsection}{Director of Training and Programs, Acceleration Consortium}{Jun 2023 - Mar 2025}{} \\
\item Delivered 5000 person hours of training content to 700 unique individuals across 20 events
\item Developed 200 hours of software and robotics course content for \href{https://ac-microcourses.readthedocs.io/}{\textit{Autonomous Systems for Discovery}}
\item Built the \href{https://ac-training-lab.readthedocs.io/}{AC Dev Lab}, a teaching, prototyping, and research facility for self-driving labs
\item Organized and delivered nine \href{https://youtu.be/IVaWl2tL06c?si=DYrhZP6uvK5srQ58}{workshops} for data science, robotics, and software development topics
\end{rSubsection}
\begin{rSubsection}{Instructor, Intro to Statistics for Engineers, University of Utah}{Aug 2021 - Dec 2021}{} \\
\item  Instructed 37 undergraduate students in probability density functions, regression, etc.
\item Conducted guided coding tutorials in Python and the Wolfram language using statistical functions
\end{rSubsection}
\begin{rSubsection}{Teaching Assistant, Intro to Materials Science, University of Utah}{Aug 2021 - Dec 2021}{} \\
\item Assisted 22 first-year graduate students in understanding materials science concepts
\end{rSubsection}
\begin{rSubsection}{Teaching Assistant, Experimental Physics Lab, BYU}{Jun 2017 - Aug 2017}{} \\
\item Assisted students in appreciating/learning physics via magnetism, electricity, and optics experiments
\end{rSubsection}
\begin{rSubsection}{Teaching Assistant, Physics Tutorial Lab, BYU}{Aug 2017 - Dec 2017}{} \\
\item Assisted students with physics problems through fundamental theory and effective questions
\end{rSubsection}
\end{rSection}

\Needspace{10\baselineskip}
\begin{rSection}{Publications}
\renewcommand{\arraystretch}{1.2}
\begin{longtable}{m{15cm} m{1.3cm}}
% 2026
{\printpublication{cluff_quantifying_2026}} & {\em \hfill 2026} \\
% 2025 preprints
{\printpublication{sayeed_honegumi_2025}} & {\em \hfill 2025} \\
{\printpublication{tang_multitask_2025}} & {\em \hfill 2025} \\
{\printpublication{baird_bayesian_2025}} & {\em \hfill 2025} \\
{\printpublication{pelkie_democratizing_2025}} & {\em \hfill 2025} \\
{\printpublication{bairdHonegumiInterfaceAccelerating2025}} & {\em \hfill 2025} \\
% 2024
{\printpublication{tomSelfDrivingLaboratoriesChemistry2024}} & {\em \hfill 2024} \\
{\printpublication{bairdMatbenchgenmetricsPythonLibrary2024}} & {\em \hfill 2024} \\
{\printpublication{loReviewLowcostSelfdriving2024}} & {\em \hfill 2024} \\
{\printpublication{sayeedNLPMeetsMaterials2024}} & {\em \hfill 2024} \\
{\printpublication{alversonGenerativeAdversarialNetworks2024}} & {\em \hfill 2024} \\
% 2023
{\printpublication{seegmillerDiscoveringChemicallyNovel2023}} & {\em \hfill 2023} \\
{\printpublication{bairdBuildingHelloWorld2023}} & {\em \hfill 2023} \\
{\printpublication{bairdCompactnessMattersImproving2023a}} & {\em \hfill 2023} \\
{\printpublication{bairdMaterialsScienceOptimization2023}} & {\em \hfill 2023} \\
{\printpublication{baird_materials_2023a}} & {\em \hfill 2023} \\
{\printpublication{sayeedStructureFeatureVectors2023}} & {\em \hfill 2023} \\
% 2022
{\printpublication{bairdWhatMinimalWorking2022}} & {\em \hfill 2022} \\
{\printpublication{bairdXtal2pngPythonPackage2022}} & {\em \hfill 2022} \\
{\printpublication{bairdHighdimensionalBayesianOptimization2022a}} & {\em \hfill 2022} \\
{\printpublication{bairdComputationallyEfficientBarycentric2022}} & {\em \hfill 2022} \\
{\printpublication{bairdDiSCoVeRMaterialsDiscovery2022}} & {\em \hfill 2022} \\
{\printpublication{bairdHighthroughputCalculationAtomic2022}} & {\em \hfill 2022} \\
{\printpublication{bairdDatadrivenMaterialsDiscovery2022}} & {\em \hfill 2022} \\
{\printpublication{titiriciSustainableMaterialsRoadmap2022}} & {\em \hfill 2022} \\
% 2021 and earlier
{\printpublication{bairdQuantitativeCartographyGrain2021}} & {\em \hfill 2021} \\
{\printpublication{bairdFiveDegreeoffreedomProperty2021}} & {\em \hfill 2021} \\
{\printpublication{fullwoodDeterminingGrainBoundary2021}} & {\em \hfill 2021} \\
{\printpublication{johnsonInferenceUncertaintyPropagation2021}} & {\em \hfill 2021} \\
{\printpublication{snowGrainBoundaryStructureproperty2021}} & {\em \hfill 2021} \\
{\printpublication{kurniawanGrainBoundaryStructure2019}} & {\em \hfill 2019} \\
{\printpublication{barrettCarbonMonolithScaffolding2017}} & {\em \hfill 2017} \\
\end{longtable}
\end{rSection}

\Needspace{10\baselineskip}
\begin{rSection}{Presentations}
\renewcommand{\arraystretch}{1.5}
\begin{longtable}{m{13cm} m{3.5cm}}

{\textit{An Autonomous Vertical Cloud Lab for Additively Manufactured Aerospace Alloys}. Materials Research Society (MRS) Fall Meeting, 2025. Oral Presentation.} & {\em \hfill Dec 2025 \par \hfill Boston, MA} \\
{\textit{SDL101: Self-Driving Lab Fundamentals and Practical Implementation}. Accelerate Conference, 2025. Training Workshop (2.5 hr).} & {\em \hfill Aug 2025 \par \hfill Toronto, ON} \\
{\textit{Autonomous Cloud Experimentation Demos at the AC Training Lab}. Acceleration Consortium Research Symposium, University of Toronto, 2024. Oral Presentation.} & {\em \hfill Oct 2024 \par \hfill Toronto, ON} \\
{\textit{Teaching AI and Automation with Physical Hardware}. Accelerate Conference, University of British Columbia, 2024. Oral Presentation.} & {\em \hfill Aug 2024 \par \hfill Toronto, ON} \\
{\textit{So, you want to build a self-driving lab?} Accelerate Conference, University of British Columbia, 2024. Oral Presentation.} & {\em \hfill Aug 2024 \par \hfill Toronto, ON} \\
{\textit{So, you want to build a self-driving lab?} AC Undergraduate Professional Development Seminars, University of Toronto, 2024. Oral Presentation.} & {\em \hfill Jun 2024 \par \hfill Toronto, ON} \\
{\textit{A Gentle Introduction to Bayesian Optimization}. Materials Science \& Engineering, North Carolina State University, 2024. Invited Guest Lecture.} & {\em \hfill Apr 2024 \par \hfill (virtual)} \\
{\textit{A Gentle Introduction to Bayesian Optimization}. Computer Science, University of Chicago, 2024. Invited Guest Lecture.} & {\em \hfill Mar 2024 \par \hfill (virtual)} \\
{\textit{Microcredentials Courses and Training Facility hosted by the Acceleration Consortium}. German-Canadian Materials Acceleration Centre (GCMAC) Summer School, 2023. Oral Presentation.} & {\em \hfill Sep 2023 \par \hfill Karlsruhe, Germany} \\
{\textit{Real-world Bayesian Optimization in Materials Science}. German-Canadian Materials Acceleration Centre Summer School, 2023. Training Workshop (1.5 hr).} & {\em \hfill Sep 2023 \par \hfill Karlsruhe, Germany} \\
{\textit{Review of Low-cost SDLs: The Frugal Twin Concept}. Accelerate Conference, 2023. Oral Presentation.} & {\em \hfill Aug 2023 \par \hfill Toronto, ON} \\
{\textit{SDL Software Best Practices}. Accelerate Conference, 2023. Training Workshop (1.5 hr).} & {\em \hfill Aug 2023 \par \hfill Toronto, ON} \\
{\textit{\href{https://youtu.be/IVaWl2tL06c?si=SsoNCR_Ny3tuPvFX}{Real-world Bayesian Optimization}}. Accelerate Conference, 2023. Training workshop (1.5 hr).} & {\em \hfill Aug 2023 \par \hfill Toronto, ON} \\
{\textit{\href{https://youtu.be/LxIiWSfkhCE?si=9wVu4S7Ar8DhtuJ0&t=163}{Data-driven Materials Discovery: Novel Discovery, Bayesian Optimization, and Self-driving Labs}}. University of Utah Dissertation Defense, 2023. Virtual.} & {\em \hfill Aug 2023 \par \hfill Virtual} \\
{\textit{Compactness Matters: Improving Bayesian Optimization Efficiency of Materials Formulations through Invariant Search Spaces}. 152nd Annual Meeting \& Exhibition of The Minerals, Metals \& Materials Society (TMS), 2023. Oral Presentation.} & {\em \hfill Mar 2023 \par \hfill San Diego, CA} \\
{\textit{What is a minimal working example for a self-driving laboratory?} 152nd Annual Meeting \& Exhibition of The Minerals, Metals \& Materials Society (TMS), 2023. Oral Presentation.} & {\em \hfill Mar 2023 \par \hfill San Diego, CA} \\
{\textit{Integrating Simulations and Experiments into a Materials Discovery Campaign: Case Studies in Multi-Fidelity Bayesian Optimization}. 152nd Annual Meeting \& Exhibition of The Minerals, Metals \& Materials Society (TMS), 2023. Oral Presentation.} & {\em \hfill Mar 2023 \par \hfill San Diego, CA} \\
{\textit{Inverse design through generative materials informatics via xtal2png}. 152nd Annual Meeting \& Exhibition of The Minerals, Metals \& Materials Society (TMS), 2023. Oral Presentation.} & {\em \hfill Mar 2023 \par \hfill San Diego, CA} \\
{\textit{How Should You Select an Algorithm for a Materials Discovery Campaign?} 152nd Annual Meeting \& Exhibition of The Minerals, Metals \& Materials Society (TMS), 2023. Oral Presentation.} & {\em \hfill Mar 2023 \par \hfill San Diego, CA} \\
{\textit{How Do You Optimize Your Parameters? Realistically Complex Hyperparameter Optimization of 23 Parameters of a Black Box Function over a Realistically Low Budget of 100 Iterations}. 152nd Annual Meeting \& Exhibition of The Minerals, Metals \& Materials Society (TMS), 2023. Poster.} & {\em \hfill Mar 2023 \par \hfill San Diego, CA} \\
{\textit{Data-driven Materials Discovery}. Materials Science \& Engineering Graduate Seminar, 2023. Oral Presentation.} & {\em \hfill Mar 2023 \par \hfill San Diego, CA} \\
{\textit{Adaptive Experimentation for Science: A More Efficient Version of Design of Experiments}. 18th IEEE International Conference on e-Science, 2022. Virtual Tutorial (3 hr).} & {\em \hfill Oct 2022 \par \hfill Salt Lake City, UT} \\
{\textit{\href{https://youtu.be/LNkRjByzeZg?si=04bLIokgA4k_HTFB&t=33}{What is a minimal working example for self-driving laboratories?}}. Accelerate Conference, 2022. Invited Talk.} & {\em \hfill Aug 2022 \par \hfill Salt Lake City, UT} \\
{\textit{Materials Properties Prediction}. Machine Learning for Materials Hard and Soft, Erwin Schrödinger International Institute for Mathematics and Physics (ESI), University of Vienna, 2022. Oral Presentation/Tutorial (1.5 hr).} & {\em \hfill Jul 2022 \par \hfill Vienna, Austria} \\
{\textit{Ensemble of State-of-the-art Property Prediction Machine Learning Algorithms}. 151st Annual Meeting \& Exhibition of The Minerals, Metals \& Materials Society (TMS), 2022. Oral Presentation.} & {\em \hfill Mar 2022 \par \hfill Anaheim, CA} \\
{\textit{5D Grain Boundary Energy Landscapes, Paths and Correlations from Bayesian Inference}. 151st Annual Meeting \& Exhibition of The Minerals, Metals \& Materials Society (TMS), 2022. Oral Presentation.} & {\em \hfill Mar 2022 \par \hfill Anaheim, CA} \\
{\textit{Bayesian Inference and Uncertainty Quantification of Grain Boundary Properties}. 150th Annual Meeting \& Exhibition of The Minerals, Metals \& Materials Society (TMS), 2021. Oral Presentation.} & {\em \hfill Mar 2021 \par \hfill Virtual} \\
{\textit{Backtracking 5DOF Grain Boundary Hydrogen Diffusivities in High-Purity Nickel: Experimentation, Localization Techniques,
and Inverse Problem Theory}. Materials Science \& Technology (MS\&T) 2019. Oral Presentation.} & {\em \hfill Oct 2019 \par \hfill Portland, OR} \\
{\textit{Technical Criteria of Commercially Viable Batteries}. 2018 Annual Meeting of International Coalition for Energy Storage and Innovation (ICESI). Oral Presentation.} & {\em \hfill Jan 2018 \par \hfill Dalian, China} \\
{\textit{Modified Graphene for Lithium-Air Battery Cathode}. NASA Langley Research Center. Advanced Materials and Processing Branch. Oral Presentation.} & {\em \hfill May 2017 \par \hfill Hampton, VA} \\
{\textit{Carbon Nanotube Interlayer for High Coulombic Efficiency Lithium-Sulfur Batteries without Unsafe Additives}. Materials Science \& Technology (MS\&T) 2016 Poster Session.} & {\em \hfill Oct 2016 \par \hfill Salt Lake City, UT} \\
{\textit{High Energy Lithium-Sulfur Batteries: An Application of Carbon Nanotubes}. Harold B. Lee Library Student Research Poster Competition.} & {\em \hfill Oct 2016 \par \hfill Provo, UT} \\
\end{longtable}
\end{rSection}

\Needspace{10\baselineskip}
\begin{rSection}{Community Contributions}
\renewcommand{\arraystretch}{1.5}
\begin{tabular}{m{13.9cm} m{2.3cm}}
{\href{https://github.com/sgbaird}{17k Code Contributions} (15k commits, 1.5k issues, 913 pull requests, 151 discussions)} & {\em \hfill 2018-2024} \\
{\href{https://github.com/sgbaird?tab=stars}{GitHub starred lists} comprising 1.7k repositories across 32 categories} & {\em \hfill 2018-2024} \\
{\href{https://github.com/AccelerationConsortium}{Creation of Acceleration Consortium GitHub Organization} (external \& internal use)} & {\em \hfill 2023-2024} \\
{Guest Editor for \textit{npj Computational Materials} collection: \textit{Self-driving Labs for Chemistry and Materials Science}} & {\em \hfill 2024-present} \\
{Journal Reviewer: \textit{npj Computational Materials}, \textit{Nature Communications}, \textit{Digital Discovery}, \textit{Computational Materials Science}, \textit{Journal of Open Source Software}} & {\em \hfill 2020-present} \\
{Review Panelist: National Science Foundation (NSF) TIP Directorate} & {\em \hfill 2026} \\
\end{tabular}
\end{rSection}

\Needspace{10\baselineskip}
\begin{rSection}{Grants}
\renewcommand{\arraystretch}{1.5}
\begin{tabular}{m{14.2cm} m{2.1cm}}
{NSF EAGER SSMCDAT2023 (\textbf{US\$200,000}): Natural Language Processing and Large Language Models for Automated Extraction of Materials Chemistry Data from Scientific Literature} & {\em \hfill 2023-2025} \\
Fehrmann Materials GmbH Research Contract (\textbf{US\$34,000}): Joint strength and ductility machine learning representations for aluminum alloys & {\em \hfill 2020} \\
BYU ORCA grant (US\$1,500): Cell-level Comparisons between Literature and Industry for Lithium-Sulfur, Lithium-Ion, Lithium-Oxygen, and other Next-Gen Batteries & {\em \hfill 2017-2018} \\
BYU Kennedy Center travel grant (US\$1,000): ICESI 2018 in Dalian, China & {\em \hfill 2018} \\
\end{tabular}
\end{rSection}

\Needspace{10\baselineskip}
\begin{rSection}{Awards}
\renewcommand{\arraystretch}{1.5}
\begin{tabular}{m{14.2cm} m{2.1cm}}
{Outstanding Ph.D. Student, Materials Science \& Engineering, University of Utah} & {\em \hfill 2023} \\
{"Call for Papers" winner --- invited talk and full travel sponsorship, Accelerate '22} & {\em \hfill 2022} \\
{Graduate Student Recognition Award, Mech. Eng., Brigham Young University} & {\em \hfill 2021} \\
{Gregory B. McKenna Fellowship, Materials Science \& Engineering, University of Utah} & {\em \hfill 2020-2023} \\
\end{tabular}
\end{rSection}

\Needspace{10\baselineskip}
\begin{rSection}{Skills}
\begin{longtable}{ l l l }
{\bf Software Tools} \\
\hline
Python & Mathematica & MATLAB \\
LaTeX & git/GitHub & Slurm \\
Command line interfaces & GitHub Actions & Unit testing (pytest) \\
Visual Studio Code & Optimization (Bayesian, genetic) & Fusion 360 \\
Ax Platform (Bayesian opt.) & Sphinx documentation & ReadtheDocs \\
Prefect (workflow orchestration) & Streamlit web apps & Gradio web apps \\
Jinja2 & Anaconda packaging & PyPI packaging \\
Hugging Face & paho-mqtt (MQTT) & MicroPython \\
Bash & Linux & Jupyter Notebooks \\
OpenCV (Computer Vision) & AprilTag (fiducial markers) & GitHub Codespaces \\
GitHub Classroom & NumPy & SciPy \\
scikit-learn & Matplotlib & Plotly \\
MongoDB Atlas (database) & Miniconda & Python debugging \\
Markdown & PyScaffold (project templates) & AWS Lambda \\
PreForm (SLA printing) & Bambu Studio (FDM printing) & Discourse \\
RayTune & Numba & PyScript \\
\\
{\bf Synthesis} & &\\
\hline
Furnace design \& operation & high-vacuum design & Photolithography \\
CVD operation \& maintenance & LPCVD operation \& maintenance & Thermal evaporation \\
Glovebox operation \& maintenance & Picosecond laser machining & Optics design \\
Etching (chemical, plasma) & Polishing (electro/mechanical) & FIB micro-machining \\
Automated liquid handling & Automated solid dosing & Six-axis robotics
\\
\\
{\bf Characterization} & & \\
\hline
Microscopy (optical, SEM) & EBSD & EDS \\
XRD & Battery cycling & Laser interferometry \\ Raman spectroscopy & Electrochemical permeation & FTIR spectroscopy \\
\end{longtable}
\end{rSection}

\end{document}

%\begin{rSection}{Research Interests}
%\end{rSection}

%\begin{rSection}{Coursework}

%\end{rSection}
